\deu{\section{Stochastik}\index{Stochastik}}
\eng{\section{Stochastics}\index{Stochastics}}
\setcounter{aufgabenNummer}{400}
  \renewcommand{\kAufgabenBuchstabe}{S}

\subsection{\deu{Grundlagen}\eng{Basics}}

%% %%%%%%%%%%%%%%%%%%%%%%% Aufgabe %%%%%%%%%%%%%%%%%%%%%%%%%%%%%%%%%%%%%%%%%%%%%%%%


\kNiveauAufgabe{\deu{Mengenschreibweise von Ereignissen}\eng{Quantity notation of events}:
\deu{Eine Urne enthält 10 Kugeln mit den Zahlen 0 bis 9. Eine Kugel wird
gezogen.}
\eng{An urn contains 10 balls with the numbers 0 to 9. One single ball is drawn.}

\begin{enumerate}[label=\alph*)]
\item
\deu{
Geben Sie folgende Ereignisse in aufzählender Form an:

A: Die Kugel trägt als Zahl eine Primzahl.\\
B: Die Zahl auf der Kugel ist durch 5 teilbar.\\
C: Die Zahl auf der Kugel ist ungerade.\\
D: Die Zahl auf der Kugel ist > 8.\\
E: Die Zahl auf der Kugel ist > 20.\\
F: Die Zahl auf der Kugel ist eine Quadratzahl.\\
G: Sicheres Ereignis
}
\eng{
State the following events in enumerative form:

A: The number on the sphere is a prime number.\\
B: The number on the sphere is divisible by 5.\\
C: The number on the sphere is odd.\\
D: The number on the sphere is > 8.\\
E: The number on the sphere is > 20.\\
F: The number on the ball is a square number.\\
G: Certain event
}%% end eng

\deu{
\item Zählen Sie alle Elementarereignisse auf.
\item Es wurde eine 7 gezogen. Welche der Ereignisse aus Teilaufgabe a) sind eingetroffen?
\item Welche Ereignispaare aus a) schliessen sich gegenseitig aus?
\item Beschreiben Sie die Gegenereignisse der unter a) aufgeführten
Ereignisse
}

\eng{
\item List all elementary events.
\item A 7 was drawn. Which of the events from subtask a) occurred?
\item Which pairs of events from a) are mutually exclusive?
\item Describe the counter-events of the events listed under a)
}
\end{enumerate}
\kVerweiseAltesKompendium{43}{5}
\kPlatzFuerBerechnungen{8}
}{%% Lösungen

a) $A = \{2, 3, 5, 7\}, B = \{0, 5\}, C = \{1, 3, 5, 7, 9\}, D = \{9\}, E = \{\}, F = \{0, 1, 4, 9\},\\ G = \{0, 1, 2, 3, 4, 5, 6, 7, 8, 9\}$

b) $\{0\}, \{1\}, \{2\}, \{3\}, \{4\}, \{5\}, \{6\}, \{7\}, \{8\}, \{9\}$

c) $A, C, G$

d) $A \cap D = \{\}, A \cap E = \{\}, A \cap F = \{\}, B \cap D = \{\}, B \cap  E = \{\},\\ C \cap E = \{\}, D \cap E = \{\}, E  \cap F = \{\}$


\kKommentar{dies ist eine Aufzählung, keine Beschreibung:}
e) $\overline{A} = \{0, 1, 4, 6, 8, 9\}, \overline{B} = \{1, 2, 3, 4, 6, 7, 8, 9\}, \overline{C} = \{0, 2, 4, 6, 8\}, \overline{D} = \{0, 1, 2, 3, 4, 5, 6, 7, 8\},\\
\overline{E} = \{0, 1, 2, 3, 4, 5, 6, 7, 8, 9\}, \overline{F} = \{2, 3, 5, 6, 7, 8\}, \overline{G} = \{\}$
}

%% %%%%%%%%%%%%%%%%%%%%%%% Aufgabe %%%%%%%%%%%%%%%%%%%%%%%%%%%%%%%%%%%%%%%%%%%%%%%%


\kNiveauAufgabe{\deu{
Der Ergebnisraum $\Omega$ und seine Teilmengen (= Ereignisse):
Aus einer grossen Produktionsserie von Leuchtdioden werden nacheinander und zufällig drei
Stück entnommen (ohne Zurücklegen). Sie werden auf «defekt» (d) bzw. «nicht defekt» (n)
hin überprüft.}
\eng{The result space $\Omega$ and its subsets (= events):
From a large production series of light-emitting diodes (LED), three pieces are taken one after the other at random (without putting them back). They are checked for «defective» (d) or «not defective» (n).
}

\begin{enumerate}[label=\alph*)]
\deu{
\item  Geben Sie den Ergebnisraum $\Omega$ an.
\item  Geben Sie die Ereignisse A bis E als Teilmengen von $\Omega$ an:
A: Die zuerst entnommene Leuchtdiode ist defekt.\\
B: Die zuerst entnommene Leuchtdiode und nur diese ist defekt.\\
C: Alle entnommenen Leuchtdioden sind defekt.\\
D: Genau zwei entnommene Leuchtdioden sind defekt.\\
E: Mindestens zwei der entnommenen Leuchtdioden sind in Ordnung.\\
}
\eng{\item Specify the result space $\Omega$.
\item Specify the events A to E as subsets of $\Omega$:

A: The first light-emitting diode removed is defective.\\
B: The first light-emitting diode removed and only this one is defective.\\
C: All removed LEDs are defective.\\
D: Exactly two removed LEDs are defective.\\
E: At least two of the removed LEDs are OK.\\
}
\end{enumerate}
\kVerweiseAltesKompendium{44}{7}
\kPlatzFuerBerechnungen{8}
}{%% Lösungen
a) $$ \Omega = \{ddd, ddn, dnd, dnn, ndd, ndn, nnd, nnn\}$$

b) $A=\{ddd, ddn, dnd, dnn\}, B = \{dnn\}, C = \{ddd\}, D = \{ddn, dnd, ndd\},\\
E = \{dnn, ndn, nnd, nnn\}$
}

%% %%%%%%%%%%%%%%%%%%%%%%% Aufgabe %%%%%%%%%%%%%%%%%%%%%%%%%%%%%%%%%%%%%%%%%%%%%%%%


\kNiveauAufgabe{
\deu{Wird beim Würfeln mit 2 gewöhnlichen Spielwürfeln Augensumme 7 (Ereignis A) oder Augensumme 8 (Ereignis B) häufiger auftreten?}
\eng{When rolling 2 ordinary dice, will the sum of 7 (event A) or the sum of 8 (event B) occur more frequently?}
\kVerweiseAltesKompendium{44}{14}
\kPlatzFuerBerechnungen{8}
}{%% Lösungen
Augensumme 7 wird häufiger auftreten.

Es gibt 6 Varianten, die Augensumme 7 zu werfen: $1+6; 2+5; 3+4; 4+3;
5+2; 6+1$

Es gibt nur 5 Varianten, die Augensumme 8 zu werfen; $2+6; 3+5; 4+4;
5+3; 6+2$

}

\newpage

\subsection{\deu{Kombinatorik}\eng{Combinatorics}}

\kNiveauAufgabe{
In einer Urne sind fünf Bälle. Ein roter (r), ein blauer (b), ein
grüner (g), ein schwarzer (s) und ein weißer (w).

Sie nehmen vier der fünf Kugeln heraus. Auf wie viele Arten ist dies
möglich, wenn

a) die Kugeln nach jedem Ziehen wieder in die Urne zurückgelegt werden und die
Reihenfolge der gezogenen Kugeln wesentlich ist?

b) die Kugeln nach jedem Ziehen wieder in die Urne zurückgelegt werden und die
Reihenfolge der gezogenen Kugeln \textbf{un}wesentlich ist?
 
c) die Kugeln nach dem Ziehen \textbf{nicht} wieder in die Urne zurückgelegt werden und die
Reihenfolge der gezogenen Kugeln wesentlich ist?

d) die Kugeln nach dem Ziehen \textbf{nicht} wieder in die Urne zurückgelegt werden und die
Reihenfolge der gezogenen Kugeln \textbf{un}wesentlich ist?

\kKommentar{Neue Aufgabe mit allen vier Fällen.}
\kPlatzFuerBerechnungen{4}
}{%% Lösungen
a) $5^4 = 625$

b) ${5+4-1 \choose 4} = 70$

c) $\frac{5!}{(5-4)!} = 120$

d) $\frac{5!}{4!} = 5$

}


%% %%%%%%%%%%%%%%%%%%%%%%% Aufgabe %%%%%%%%%%%%%%%%%%%%%%%%%%%%%%%%%%%%%%%%%%%%%%%%

\kNiveauAufgabe{
\deu{Ein Schüler will in den Sommerferien seine Schulbücher – ein Englischbuch, drei Deutschbücher, zwei Französischbücher und vier Mathematikbücher (alle Bücher sind verschieden!)
– ordnen. Sie sollen auf ein Regal gestellt werden, wobei Bücher des gleichen Fachgebietes
nebeneinander stehen sollen. Wie viele verschiedene Gruppierungsmöglichkeiten der Bücher
gibt es?}
\eng{During the summer vacation, a pupil wants to bring order into his
school books --- one English book, three German books, two French books and four math books (all books are different!).
They are to be placed on a shelf, with books from the same subject stand next to each other.
In how many different ways can the books be grouped?
}
\kVerweiseAltesKompendium{46}{15}
\kPlatzFuerBerechnungen{8}
}{%% Lösungen
 $$m = (1! \cdot{} 3! \cdot{} 2!  \cdot{} 4!) \cdot{} (4!) = 6912$$
}


%% %%%%%%%%%%%%%%%%%%%%%%% Aufgabe %%%%%%%%%%%%%%%%%%%%%%%%%%%%%%%%%%%%%%%%%%%%%%%%
\kNiveauAufgabe{
Die sieben Kollega der Dorfschule (drei Schülerinnen und vier Schüler)
gehen zusammen ins Dorfkino. Die vorderste Reihe hat genau sieben
Plätze.

Auf wie viele Arten können sich die sieben Kollega auf die sieben
Plätze anordnen?

\kPlatzFuerBerechnungen{8}
}{%% Lösung
$7! = 5040$
}%% end Aufgabe

 
%% %%%%%%%%%%%%%%%%%%%%%%% Aufgabe %%%%%%%%%%%%%%%%%%%%%%%%%%%%%%%%%%%%%%%%%%%%%%%%

\kNiveauAufgabe{\deu{Vier Personen steigen in ein Zugsabteil. Der Zug
fährt mit diesen vier Personen los und hält anschliessend noch an
sechs weiteren Haltestellen.}
\eng{Four people board a train compartment. The train departs with these four
people and then stops at six more stops.}
\begin{enumerate}[label=\alph*)]
\item \deu{Wie viele verschiedene Möglichkeiten haben die vier Personen, um den Zug zu verlassen
(es dürfen auch mehrere Personen gleichzeitig den Zug verlassen)?}
\eng{How many different ways do the four people have to get off the train

(several people can leave the train at the same time)?}
\item \deu{Wie viele Möglichkeiten gibt es, wenn jede Person an einer
andern Station aussteigt?}
\eng{How many possibilities are there if each «passenger» gets off at a
get off at a different station?}
\end{enumerate}
\kVerweiseAltesKompendium{46}{16}
\kPlatzFuerBerechnungen{8}
}{%% Lösungen
a) $6^4 = 1296$

b) $6\cdot{}5\cdot{}4\cdot{}3 = 360$
}

%% %%%%%%%%%%%%%%%%%%%%%%% Aufgabe %%%%%%%%%%%%%%%%%%%%%%%%%%%%%%%%%%%%%%%%%%%%%%%%
\kNiveauAufgabe{Ein Zahlenschloss hat fünf Ringe und jeder Ring hat
acht Ziffern.

a) Wie viele Variationen sind damit möglich?

b) Wie viele Ringe würden benötigt, wenn 1\,000\,000 Variationen
möglich sein sollten.

\kKommentar{neue Aufgabe}
}{%% Lösungen
a) $8^5 = 32\,768$
b) $8^x = 1\,000\,000 \Longrightarrow x=\log_8(10^6) \approx
6.6 \Longrightarrow \textbf{sieben} Ringe $
}%% end Aufgabe

%% %%%%%%%%%%%%%%%%%%%%%%% Aufgabe %%%%%%%%%%%%%%%%%%%%%%%%%%%%%%%%%%%%%%%%%%%%%%%%

\kNiveauAufgabe{\deu{Die Buchstaben aus dem Wort MISSISSIPPI werden durcheinandergemischt und zufällig
wieder zusammengesetzt.

a) Wieviele verschiedene Buchstabenfolgen mit und ohne Bedeutung
können dabei entstehen?}
\eng{The letters from the word MISSISSIPPI are mixed up and are then randomly reassembled.


a) How many different sequences of letters with and without meaning can
result from this procedure?}

b) \deu{Wie viele Wörter sind möglich, wenn die beiden P nebeneinander
stehen müssen?}\eng{How many possibilities are there, when the letters
P have to stand right next to each other?}

c) \deu{Wie viele, wenn die P's nicht nebeneinander stehen
dürfen?}\eng{How many, when the P's must not touch each other.}
\kKommentar{neue Aufgabe}
\kPlatzFuerBerechnungen{8}
}{%% Lösungen
a) 
$$m=\frac{11!}{2!\cdot{} 4! \cdot{} 4!} = 34\,650$$

b) \deu{Betrachte PP als einen Buchstaben}\eng{}
$$m=\frac{10!}{4! \cdot{} 4!} = 6300$$

c) $34\,650 - 6300 = 28350$
}



%% %%%%%%%%%%%%%%%%%%%%%%% Aufgabe %%%%%%%%%%%%%%%%%%%%%%%%%%%%%%%%%%%%%%%%%%%%%%%%



\kNiveauAufgabe{
\eng{translate}
\kKommentar{neue Aufgabe\\}

Die «Glorreichen Sieben» fahren nach ihrer Film-Premiere nach
Hause. Leider hat das Taxi nur Platz für vier der sieben. Drei der
glorreichen sieben gehen daher zu Fuss.

Wie viele Kombinationen sind so möglich, wer im Taxi mitfahren darf?
}{%%Lösung

$${7 \choose 4} = 35$$
}


%% %%%%%%%%%%%%%%%%%%%%%%% Aufgabe %%%%%%%%%%%%%%%%%%%%%%%%%%%%%%%%%%%%%%%%%%%%%%%%

\kNiveauAufgabe{
\deu{Lou soll für ein Fest 12 Getränkeflaschen kaufen; genau so viele wie
im Harass (= Getränkekiste) Platz finden. Es stehen 5 verschiedene Getränke zur
Auswahl. Auf wie viele Arten kann Lou den Harass füllen? Die
Reihenfolge der Flaschen im Harass spielt dabei keine Rolle.}
\eng{Lou is supposed to buy 12 beverage bottles for a party,
exactly the number that fits in a crate. There are 5 different drinks
to choose from.
In how many ways can Lou fill the crate? The order of the bottles in the crate does not matter.}
\kKommentar{neue Aufgabe}
\kPlatzFuerBerechnungen{8}
}{%%\Lsung
\eng{There are 1820 variants.}\deu{Es gibt 1820 Varianten.}}

\newpage